\documentclass[a4paper]{article}

\usepackage[utf8]{inputenc}

\usepackage{minted}
\usepackage{tikz}
\usepackage{tcolorbox}
\usepackage{amsmath}
\usepackage{amsthm}
\usepackage{amsfonts}	
\usepackage{float}
\usepackage{listings}
\usepackage{fancyhdr}


\author{Lazar Nagulov}
\title{testA}
\date{}

\begin{document}
    \include{listings}
    \maketitle
    
        \section{Templates}
    \par Templates are blueprints for creating reusable data structures or 
    types that are used in a "generate" expression. Each template has 
    a name and a set of fields, where each field is a pair consisting 
    of a string (the field name) and an expression (which defines the 
    field's value). By using data types and enumerations within these 
    expressions, you can add randomness to the values generated by the 
    template.
    \par An example can been seen in listing \ref{listing:1}.
    \begin{listing}[h]
        \begin{minted}[escapeinside=||, xleftmargin=1em, breaklines]{cpp}
template <name> {
    <name1> = <expr>;
    <name1> = <expr>;
    ...
}

template User {
    id = int;
    name = |\textcolor{purple}{string}|;
    password = |\textcolor{purple}{string}|;
    age = int;
    balance = float;
}
        \end{minted}
        \caption{Templates in testA}
        \centering
        \label{listing:1}
    \end{listing}
    \pagebreak
    \section{Enumerations}
    \par The enumerations in testA consist of a list of names and their corresponding weights.
    Weights represent the likelihood of randomly generating an enumeration variant. They are always integers, and the default weight is 1 if not specified.
    \par The enumeration declaration is shown in listing \ref{listing:2}.
    \begin{listing}[h]
        \begin{minted}[escapeinside=||, xleftmargin=1em, breaklines]{cpp}
enum <name> { 
     <variant1> [=> <expr:int>]; 
     <variant2> [=> <expr:int>]; 
     ... 
}  

enum Role {
    User => 8;
    Admin; // Default weight is 1
    Developer => 3;
}
        \end{minted}
        \caption{Enumerations in testA}
        \centering
        \label{listing:2}
    \end{listing}
    \par As shown in listng \ref{listing:3}, enumerations can be used as types in templates. Their evaluated values are always represented as strings
    \begin{listing}[h]
        \begin{minted}[escapeinside=||, xleftmargin=1em, breaklines]{cpp}
template User {
    id = int;
    name = |\textcolor{purple}{string}|;
    role = |\textcolor{purple}{Role}|;
}
        \end{minted}
        \caption{Template with enumeration field}
        \centering
        \label{listing:3}
    \end{listing}

    \section{Resources}
    \section{Generate}

\end{document}